% !TEX root = ../../report.tex

\section{Ролевая модель}

Ролевая модель используется для разграничения доступа к таблицам базы данных на уровне СУБД. Каждой роли назначается набор привилегий, определяющий допустимые операции.

В аналитической части были определены пять ролей пользователей системы. На уровне СУБД проектируются четыре роли, соответствующие авторизованным пользователям. Роль гостя не выделяется как отдельная роль СУБД, поскольку действия гостя ограничиваются просмотром общедоступных данных и авторизацией. Просмотр гостем общедоступных данных соответствует чтению каталога, карточек продавцов, отзывов и истории цен.

Описание привилегий каждой роли приведено ниже.

\textbf{marketplace\_buyer} --- роль покупателя:
\begin{itemize}
    \item чтение таблиц products, sellers, categories, product\_categories, reviews, product\_price\_history --- просмотр каталога, карточек товаров, продавцов, отзывов и истории цен;
    \item управление таблицами addresses, reviews, orders, order\_items --- работа с адресами, отзывами, заказами и позициями заказа.
\end{itemize}

\textbf{marketplace\_seller} --- роль продавца:
\begin{itemize}
    \item управление таблицами sellers, products, product\_categories --- работа с профилем продавца, товарами и их категориями;
    \item чтение таблиц categories, reviews, product\_price\_history --- просмотр категорий, отзывов и истории цен;
    \item чтение таблиц orders и order\_items, обновление таблицы orders --- просмотр заказов продавца и фиксация статусов отгрузки и доставки.
\end{itemize}

\textbf{marketplace\_admin} --- роль администратора:
\begin{itemize}
    \item управление всеми таблицами --- полный доступ к данным.
\end{itemize}

\textbf{marketplace\_analyst} --- роль аналитика:
\begin{itemize}
    \item чтение всех таблиц --- доступ к данным без возможности модификации.
\end{itemize}

Сводная таблица привилегий ролей приведена в таблице~\ref{tbl:role_matrix}.

\begin{table}[H]
    \caption{Матрица привилегий ролей}
    \label{tbl:role_matrix}
    \begin{tabular}{|l|c|c|c|c|}
        \hline
        \textbf{Таблица} & \textbf{buyer} & \textbf{seller} & \textbf{admin} & \textbf{analyst} \\ \hline
        users & --- & --- & CRUD & R \\ \hline
        sellers & R & CRUD & CRUD & R \\ \hline
        products & R & CRUD & CRUD & R \\ \hline
        categories & R & R & CRUD & R \\ \hline
        product\_categories & R & CRUD & CRUD & R \\ \hline
        addresses & CRUD & --- & CRUD & R \\ \hline
        orders & CRUD & RU & CRUD & R \\ \hline
        order\_items & CRUD & R & CRUD & R \\ \hline
        reviews & CRUD & R & CRUD & R \\ \hline
        product\_price\_history & R & R & CRUD & R \\ \hline
    \end{tabular}
\end{table}

Обозначения в таблице: C --- CREATE (INSERT), R --- READ (SELECT), U --- UPDATE, D --- DELETE. Прочерк означает отсутствие доступа к данной таблице.

Матрица описывает привилегии на уровне таблиц. Правила принадлежности отдельных записей соответствуют ограничениям предметной области: продавец работает со своим профилем, своими товарами, категориями своих товаров и заказами, покупатель --- со своими заказами, позициями заказа, адресами и отзывами. История цен формируется триггером; покупатель и продавец имеют доступ только на чтение этой таблицы.

Полный доступ администратора в матрице относится к уровню СУБД. Пользовательские сценарии администратора остаются заданными в аналитической части.
